\documentclass{book}
\usepackage{pgfplots}
\usepgfplotslibrary{colorbrewer}
\pgfplotsset{compat=1.18}

% =========================================================================
% 1. MASTER BOOK CONFIGURATION SWITCH
% =========================================================================
\newif\ifcolorbook
\colorbooktrue % <--- Set to 'true' for color/digital, 'false' for print B&W

% =========================================================================
% 2. THE UNIVERSAL BOOK-WIDE LAYOUT ENGINE
% =========================================================================
\ifcolorbook
    % --- Color Palette (Using ColorBrewer Dark2 & Viridis) ---
    \pgfplotscreateplotcyclelist{booklines}{
        {Dark2-A, thick, no markers}, {Dark2-B, thick, no markers},
        {Dark2-C, thick, no markers}, {Dark2-D, thick, no markers}
    }
    \pgfplotsset{
        masterbookstyle/.style={
            cycle list name=booklines,
            colormap/viridis,
            grid style={solid, gray!15}, % Soft grids for digital viewing
        },
        outer3D/.style={shader=faceted, faceted color=none, opacity=0.45},
        inner3D/.style={shader=faceted, faceted color=none, fill=Dark2-D}
    }
\else
    % --- Monochrome Palette (Highly distinct line styles & shading) ---
    \pgfplotscreateplotcyclelist{booklines}{
        {black,       thick, solid, no markers},
        {black!60,    thick, dashed, no markers},
        {black,       thick, dotted, no markers},
        {black!70,    thick, dashdotted, no markers}
    }
    \pgfplotsset{
        masterbookstyle/.style={
            cycle list name=booklines,
            colormap/blackwhite,
            grid style={dashed, black!20}, % Highly readable dashed grid for print
        },
        outer3D/.style={shader=faceted, faceted color=black!25, thin, opacity=0.20},
        inner3D/.style={shader=faceted, faceted color=black, thin, fill=black!75}
    }
\fi

% =========================================================================
% 3. FORCING GLOBALLY BALANCED GRAPHICS (No style collisions)
% =========================================================================
\pgfplotsset{
    % Instantly hooks into every single \begin{axis} automatically
    every axis/.append style={
        masterbookstyle,
        % Core Typography & Positioning
        tick label style={font=\small}, 
        label style={font=\small},
        % Split keys to prevent TikZ / PGFPlots compilation errors
       legend cell align=left,              
        legend style={font=\small},          
        % Clean academic layout aesthetics 
        axis lines=left,            % Academic standard (L-shaped axis)
        grid=major,                 % Uniform background grids book-wide
        enlargelimits=false,        % Prevents data hanging past axis lines
    }
}

\begin{document}

% -------------------------------------------------------------------------
% 2D Line Plot Example
% -------------------------------------------------------------------------
\begin{figure}[htbp]
\centering
\begin{tikzpicture}
\begin{axis}[
    xlabel=$x$, 
    ylabel=$y$
]
    \addplot {x};
    \addplot {x^2};
    \addplot {x^3};

\end{axis}
\end{tikzpicture}
\caption{Linear versus power rates of change.}
\end{figure}

% -------------------------------------------------------------------------
% 3D Transparent Surface Example (Seeing an inner object through an outer shell)
% -------------------------------------------------------------------------
\begin{figure}[htbp]
\centering
\begin{tikzpicture}
\begin{axis}[
    view={60}{30}, 
    xlabel=$X$, 
    ylabel=$Y$, 
    zlabel=$Z$,
    title={Universal 3D Style}
]
    % Inner object must be drawn first in the code layer
    \addplot3[
        surf, 
        inner3D, 
        domain=-0.6:0.6, 
        samples=15
    ] {0.8 * exp(-5*(x^2+y^2))};

    % Outer transparent object drawn second to overlay smoothly
    \addplot3[
        surf, 
        outer3D, 
        domain=-2:2, 
        samples=25
    ] {exp(-0.5*(x^2+y^2))};
\end{axis}
\end{tikzpicture}
\caption{Visualization of an intersecting vector landscape.}
\end{figure}

\end{document}

